\documentclass{handout}

\provideboolean{examVersion}\setboolean{examVersion}{true}\newcommand{\ifExam}{\ifthenelse{\boolean{examVersion}}}

\setboolean{examVersion}{false}
%\setboolean{examVersion}{true}

\usepackage{subfigure}

\usepackage{handoutSetup}\usepackage{handoutShortcuts} 

 
 


\begin{document}
\handoutHeader

\begin{verbatimwrite}{\jobname.title}
Consumption with Optimal Portfolio Choice
\end{verbatimwrite}
\medskip\medskip\medskip

\handoutNameMake

\setboolean{AnswersTF}{true}
\setboolean{AnswersTF}{false}


\begin{verbatimwrite}{./body.tex}
\ifExam{{\bf Consumption with Optimal Portfolio Choice.~~}}{}\handoutC{CRRA-RateRisk} shows that for a 
\cite{merton:restat}-\cite{samuelson:portfolio} consumer facing return
$\log \RiskyAlt_{t+1} \sim \mathcal{N}(\riskyAlt -
\sigma^{2}_{\riskyAlt}/2,\sigma^{2}_{\riskyAlt})$ on the only financial
asset available, the optimal marginal propensity
to consume is approximately
\begin{equation}\begin{gathered}\begin{aligned}
        \MPC & \approx  {\riskyAlt}-\CRRA^{-1}(\riskyAlt-\DiscRate) - (\CRRA-1)\left(\sigma_{\riskyAlt}^{2}/2\right) \label{eq:MPC}
\end{aligned}\end{gathered}\end{equation}
\ifExam{
\begin{enumerate}
\item Use this equation to discuss the parametric restriction(s) under which an increase in 
unavoidable financial risk $\sigma_{\riskyAlt}^{2}/2$ will cause a (precautionary) decline in consumption (the risk is unavoidable because we have assumed the risky asset is the only financial asset available).
Explain whether, under such parameter value(s), the income effect of an increase in $\riskyAlt$ outweighs
the substitution effect, {\it and explain in words the intution for both the income effect and the substitution effect.}
\answer{The restriction is $\CRRA>1$, which guarantees that an increase in $\sigma_{\riskyAlt}^{2}$ reduces $\MPC$.  Under this restriction, the income effect outweighs the substitution effect.  For a discussion of the two effects, see \handoutC{PerfForesightCRRA}.}{\vspace{0.1in}} % end answer
}{where the precautionary effect of financial risk on the MPC is
captured by the $-(\CRRA-1)\left(\sigma_{\riskyAlt}^{2}/2\right)$ term.  
Since $\CRRA > 1$ by assumption, this equation yields the plausible 
conclusion that an increase in {\it unavoidable} financial risk $\sigma_{\riskyAlt}^{2}$ reduces
the level of consumption.} % end \ifExam

\renewenvironment{quote}{%
  \list{}{%
    \leftmargin-0.75cm   % this is the adjusting screw
    \rightmargin\leftmargin
  }
  \item\relax
}
{\endlist}


\ifExam{
\begin{quote}
For the remainder of the question, assume that the consumer}{We are interested here in understanding how the results change when the consumer} can choose how much to invest in the risky asset, so that financial risk can be avoided by reducing the share of the portfolio allocated to the risky asset.  \handoutA{Portfolio-CRRA} derives the portfolio share $\riskyshare$
  that an optimizing consumer will invest in a risky asset earning
  return $\log \RiskyAlt_{t+1} \sim \mathcal{N}(\riskyAlt -
  \sigma^{2}_{\riskyAlt}/2,\sigma^{2}_{\riskyAlt})$ -- so that $\log \RiskyAlt =
  \riskyAlt$ (where the subscriptless version of a variable denotes
  its expectation unless otherwise noted, e.g.\ $\RiskyAlt \equiv \Ex_{t}[\RiskyAlt_{t+1}]$.\footnote{\ifExam{}{(See
  \LogMeanMPS, and note that $\RiskyAlt=\Risky$ under the given assumptions; that is, an increase in the degree of risk $\sigma^{2}_{\riskyAlt}$ does not change the expected return factor in levels).}. Note that the subscriptless version of the log of the risky return, $\riskyAlt_{t+1}$, is {\it not} equal to its mean: $\Ex_{t}[\riskyAlt_{t+1}] = \riskyAlt-\sigma^{2}_{\riskyAlt}/2$.}  The remaining proportion $(1-\riskyshare)$
  of the portfolio earns a riskless return $\rfree = \log \Rfree$,\footnote{The bold
    font is used for the risky return and the narrow font for the safe
    return.} and we write the log  expected return premium factor as
  $\EPrem \equiv \RiskyAlt/\Rfree$ with log expected return premium $\EpremLog
  \equiv \log \EPrem = \riskyAlt-\rfree = \log \RiskyAlt/\Rfree$; 
  optimal choice of $\riskyshare$ yields a portfolio whose realized return factor is written
$\Rport$
\ifExam{}{and the log of whose realization is well approximated by\footnote{See \handoutA{Portfolio-CRRA} for a figure examining the quality of the approximation.}
\begin{equation}\begin{gathered}\begin{aligned} 
    \rport_{t+1}-\rfree & \approx  \riskyshare \overbrace{(\riskyAlt_{t+1}-\rfree)}^{\equiv \EpremLog_{t+1}} + \riskyshare(1-\riskyshare)\sigma^{2}_{\riskyAlt}/2 
\\   & =   \riskyshare (\riskyAlt_{t+1}-\rfree+\sigma^{2}_{\riskyAlt}/2) - \riskyshare^{2} \sigma^{2}_{\riskyAlt}/2, \label{eq:portLogPrem}
\end{aligned}\end{gathered}\end{equation}
whose variance (using \SumNormsIsNorm~and \NormTimes) is
\begin{equation}\begin{gathered}\begin{aligned}
\sigma^{2}_{\rport} & =  \riskyshare^{2} \sigma^{2}_{\riskyAlt},   \label{rportVar}
\end{aligned}\end{gathered}\end{equation}
and since $\Ex_{t}[\riskyAlt_{t+1}]=\riskyAlt-\sigma^{2}_{\riskyAlt}/2$ the expectation of \eqref{eq:portLogPrem} 
will be
\begin{equation}\begin{gathered}\begin{aligned} 
    \Ex_{t}[\rport_{t+1}-\rfree] & \approx  (\riskyAlt- \sigma^{2}_{\riskyAlt}/2 - \rfree + \sigma^{2}_{\riskyAlt}/2)\riskyshare-\riskyshare^{2}\sigma^{2}_{\riskyAlt}/2
\\   & =   \EpremLog \riskyshare -\riskyshare^{2}\sigma^{2}_{\riskyAlt}/2
\end{aligned}\end{gathered}\end{equation}
so that, using \SumNormsIsNorm,
\begin{equation}\begin{gathered}\begin{aligned}
  \rport_{t+1}-\rfree & \sim  \mathcal{N}(\EpremLog \riskyshare -\riskyshare^{2}\sigma^{2}_{\riskyAlt}/2,\riskyshare^{2}\sigma^{2}_{\riskyAlt}) \label{eq:rportDist}
\end{aligned}\end{gathered}\end{equation}
yielding, using \LogELogNormTimes, }
\begin{equation}\begin{gathered}\begin{aligned}
 \log \Rport \equiv \log \Ex_{t}[\Rport_{t+1}] & =  \riskyshare \EpremLog + \rfree \label{LogERport},
\end{aligned}\end{gathered}\end{equation}
and \handoutA{Portfolio-CRRA} shows that under these circumstances the optimal risky
portfolio share is well approximated by
\begin{equation}\begin{gathered}\begin{aligned}
         \riskyshare  & \approx  \left(\frac{\EpremLog}{\CRRA \sigma^{2}_{\riskyAlt}}\right)
 \label{eq:riskyshareMS}
\end{aligned}\end{gathered}\end{equation}
\ifExam{while the variance of the return on the optimally-chosen risky portfolio is approximately
\begin{equation}\begin{gathered}\begin{aligned}
\sigma^{2}_{\rport} & =  \riskyshare^{2} \sigma^{2}_{\riskyAlt}.   \label{rportVar}
\end{aligned}\end{gathered}\end{equation}
\end{quote}}{} % end ifExam

\renewenvironment{quote}{%
  \list{}{%
    \leftmargin+0.75cm   % this is the adjusting screw
    \rightmargin\leftmargin
  }
  \item\relax
}
{\endlist}


\newcommand{\ExplainWhyRiskShrinksC}{The way to understand this is to break down the response to the
increase in riskiness into two components. The first is the direct
precautionary effect examined in \handoutM{CRRA-RateRisk}, which works
in the way intuition suggests (more risk implies lower consumption);
the second effect is the `portfolio rebalancing' effect, examined in
\handoutA{Portfolio-CRRA}: An increase in riskiness makes the consumer
cchoose to invest less in the risky asset.  What \eqref{eq:PrecEffectRaw}
tells us is that the consumer's `flight from risk' is so effective in
reducing riskiness (because the portfolio share $\riskyshare$ enters
as a {\it squared} term in \eqref{eq:PrecEffectRaw}) that the
riskiness of the {\it portfolio} actually {\it declines} as the riskiness of
the risky asset increases.  (This result was also derived in
\handoutM{Portfolio-CRRA}).  Now the overall reduction in consumption
coming through the precautionary term makes sense: Precautionary
saving is less than before {\it because the consumer optimally chooses less
  exposure to risk} than before.}


\ifExam{
\item Show that the precautionary effect of rate-of-return risk on the precautionary contribution to the MPC, 
after taking account of optimal portfolio adjustment, is
\begin{equation}
 -(\CRRA-1)\left(\frac{(\EpremLog/\CRRA)^{2}}{2 \sigma_{\riskyAlt}^{2}}\right)  \label{eq:PrecEffect}.
\end{equation}
and explain the intuition for the result that the size of the precautionary effect {\it shrinks} as the risk grows larger.
\answer{
\noindent so that substituting from \eqref{eq:riskyshareMS}, the precautionary effect on consumption after taking account of optimal portfolio adjustment is
\begin{equation}\begin{gathered}\begin{aligned}
 -(\CRRA-1)\left(\sigma_{\rport}^{2}/2\right) & =  -(\CRRA-1)\left(\frac{\overbrace{\riskyshare^{2}\sigma_{\riskyAlt}^{2}}^{\text{from \eqref{rportVar}}}}{2}\right) 
\\ & =  -(\CRRA-1)\left(\frac{(\EpremLog/\CRRA)^{2}}{2 \sigma_{\riskyAlt}^{2}}\right)  
\end{aligned}\end{gathered}\end{equation}

\ExplainWhyRiskShrinksC
}{} % end answer
}{\noindent so that substituting from \eqref{eq:riskyshareMS}, the precautionary effect after taking account of optimal portfolio adjustment is
\begin{equation}\begin{gathered}\begin{aligned}
 -(\CRRA-1)\left(\sigma_{\rport}^{2}/2\right) & =  -(\CRRA-1)\left(\frac{\overbrace{\riskyshare^{2}\sigma_{\riskyAlt}^{2}}^{\text{from \eqref{rportVar}}}}{2}\right) \label{eq:PrecEffectRaw}
\\ & =  -(\CRRA-1)\left(\frac{(\EpremLog/\CRRA)^{2}}{2 \sigma_{\riskyAlt}^{2}}\right) 
\end{aligned}\end{gathered}\end{equation}
which says that (for $\CRRA>1$) the absolute size of the precautionary term \textit{shrinks} as the risk grows larger.
} % end ifExam

\ifExam{}{
To put it another way, when $\CRRA>1$ the effect of the precautionary
motive in reducing consumption gets {\it smaller} as the risk gets
larger (just to confirm that you read that right, I'll say it a third
way: an {\it increase} in the riskiness of the risky asset causes
consumption to {\it rise}).  At first, this seems bizarre: Intuition
suggests that in reality that people {\it cut back} on consumption in
the face of greater risk (financial or nonfinancial).

\ExplainWhyRiskShrinksC
} % end ifExam 


\ifExam{
\item Now use the foregoing results to show that the MPC $\MPC$ can be rewritten in terms of 
exogenous parameters (including the riskfree interest factor $\rfree$) as 
\begin{equation}\begin{gathered}\begin{aligned}
        \MPC & \approx  \rfree-\CRRA^{-1}(\rfree-\DiscRate)+ (\CRRA-1)\left(\frac{(\EpremLog/\CRRA)^{2}}{2 \sigma_{\riskyAlt}^{2}}\right)
\end{aligned}\end{gathered}\end{equation}
and use this equation to explain why the total effect of an increase in risk on consumption is {\it positive} and {\it explain why this result comes about.}
\answer{See \handoutA{C-With-Optimal-Portfolio} for the result and the logic.}{}
}{
The fact that precautionary saving diminishes when risk is larger does
not, by itself, tell us whether consumption increases in response to
an increase in risk; in addition to the precautionary channel, the MPC
is also affected by the decline in the portfolio's expected rate of return
that results from the consumer's choice to invest less in the
high-expected-return risky asset and more in the low-return safe asset.
Substituting \eqref{eq:riskyshareMS} into \eqref{LogERport} indicates
that the log of the expected portfolio excess return factor becomes
\begin{equation}\begin{gathered}\begin{aligned}
 \log \Rport/\Rfree \equiv \log \Ex_{t}[\Rport_{t+1}/\Rfree] & =  \frac{\EpremLog^{2}}{\CRRA \sigma^{2}_{\riskyAlt}}
\end{aligned}\end{gathered}\end{equation}
so an increase in $\sigma^{2}_{\riskyAlt}$ can have quite a powerful effect in reducing the consumer's
expected portfolio return.

As with the change in the riskfree rate analyzed in \handoutC{PerfForesightCRRA}, this change in the portfolio
return has two consequences: An income and a substitution effect,
captured respectively by the first and second terms of \eqref{eq:MPC}.
Substituting $\rport = \EpremLog^{2}/\CRRA \sigma^{2}_{\riskyAlt}+\rfree$ and $\sigma_{\rport}^{2}=(\EpremLog/\CRRA)^2/\sigma_{\riskyAlt}^{2}$ into \eqref{eq:MPC} yields
\begin{equation}\begin{gathered}\begin{aligned}
        \MPC & \approx  {\rport}-\CRRA^{-1}(\rport-\DiscRate) - (\CRRA-1)\left(\sigma_{\rport}^{2}/2\right) \label{eq:MPCsub}
\\      & =  \CRRA^{-1}\DiscRate+{\rport}(1-\CRRA^{-1}) - (\CRRA-1)\left(\frac{(\EpremLog/\CRRA)^{2}}{2 \sigma_{\riskyAlt}^{2}}\right)
\\      & =  \CRRA^{-1}\DiscRate+(\rfree+\EpremLog^{2}/\CRRA \sigma^{2}_{\riskyAlt})(1-\CRRA^{-1}) - (\CRRA-1)\left(\frac{(\EpremLog/\CRRA)^{2}}{2 \sigma_{\riskyAlt}^{2}}\right)
\\      & =  \rfree-\CRRA^{-1}(\rfree-\DiscRate)+\underbrace{(\CRRA-1)(\EpremLog^{2}/\CRRA^{2} \sigma^{2}_{\riskyAlt})}_{\text{net inc and sub effect}} - (\CRRA-1)\left(\frac{(\EpremLog/\CRRA)^{2}}{2 \sigma_{\riskyAlt}^{2}}\right)
.
%\\      & =  \rfree-\CRRA^{-1}(\rfree-\DiscRate)+ (\CRRA-1)\left(\frac{(\EpremLog/\CRRA)^{2}}{2 \sigma_{\riskyAlt}^{2}}\right)
\end{aligned}\end{gathered}\end{equation}

Thus, for $\CRRA > 1$ the income effect outweighs the substitution
effect so that when an increase in financial risk causes consumers to
shy away from the risky asset the reduction in consumption from the
drop in income (`the income effect') is larger than the increase in
consumption from the lowering of the incentive to delay consumption
(`the substitution effect').

Note, finally, that the net of the income and substitution effects is (remarkably) of exactly the same functional
form as the precautionary effect, but of opposite sign and twice as large.  So the combination of all 
three effects (income, substitution, precautionary) yields:
\begin{equation}\begin{gathered}\begin{aligned}
        \MPC & \approx  \rfree-\CRRA^{-1}(\rfree-\DiscRate)+ (\CRRA-1)\left(\frac{(\EpremLog/\CRRA)^{2}}{2 \sigma_{\riskyAlt}^{2}}\right)
\end{aligned}\end{gathered}\end{equation}
and we can see that the overall effect of the increase in financial
risk is to reduce the marginal propensity to consume because the net
of the income and substitution effects channels outweighs the precautionary effect (see
\handoutM{PerfForesightCRRA} for a refresher on income and
substitution effects, as well as human wealth effects (which are
absent here but could easily be added by giving the consumer an
expected stream of future noncapital income)).

Figure~\ref{fig:MPCvsSigma} confirms, both using the approximate formulas
derived above and using a numerically exact solution, that the MPC declines
as risk increases.  

} % end ifExam



\ifExam{
\item Now show that the effect of an increase in risk on saving is given by 
\begin{equation}\begin{gathered}\begin{aligned}
  \rport - \MPC & =  \CRRA^{-1}(\rfree-\DiscRate) + \left(\frac{\CRRA+1}{2}\right)\left(\frac{\EpremLog^{2}}{\sigma^{2}_{\riskyAlt}\CRRA^{2}}\right)
\end{aligned}\end{gathered}\end{equation}
and {\it explain how this result relates to the results obtained earlier}.  Comment, in particular, on the role 
played by our assumption that the consumer has no labor income (Hint: What further effect of interest rates is omitted when the consumer has no labor income?)
\answer{See the discussion in \handoutA{C-With-Optimal-Portfolio}.}{}
\end{enumerate}
}{
Finally, in addition to the effects on consumption, we are interested in the 
effects on saving.  Since saving is income minus consumption, this means
we need to know the effects on income.  But income for our consumer is 
entirely from his portfolio investments, so the income effect is captured
simply by $\rport$.  The `saving effect' is therefore captured by
\begin{equation}\begin{gathered}\begin{aligned}
  \rport - \MPC & \approx  \CRRA^{-1}(\rport-\DiscRate)+(\CRRA-1)\left(\frac{(\EpremLog/\CRRA)^{2}}{2 \sigma_{\riskyAlt}^{2}}\right)
\\ & =  \CRRA^{-1}\left(\frac{\EpremLog^{2}}{\CRRA \sigma^{2}_{\riskyAlt}}+\rfree-\DiscRate\right)+(\CRRA-1)\left(\frac{(\EpremLog/\CRRA)^{2}}{2 \sigma_{\riskyAlt}^{2}}\right)
\\ & =  \CRRA^{-1}(\rfree-\DiscRate) + \frac{\EpremLog^{2}}{\CRRA^{2} \sigma^{2}_{\riskyAlt}}+(\CRRA-1)\left(\frac{(\EpremLog/\CRRA)^{2}}{2 \sigma_{\riskyAlt}^{2}}\right)
\\ & =  \CRRA^{-1}(\rfree-\DiscRate) + \frac{\EpremLog^{2}/\CRRA^{2}}{ \sigma^{2}_{\riskyAlt}}+(\CRRA-1)\left(\frac{(\EpremLog/\CRRA)^{2}}{2 \sigma_{\riskyAlt}^{2}}\right)
\\ & =  \CRRA^{-1}(\rfree-\DiscRate) + \left(\frac{\CRRA+1}{2}\right)\left(\frac{\EpremLog^{2}}{\sigma^{2}_{\riskyAlt}\CRRA^{2}}\right)
\end{aligned}\end{gathered}\end{equation}
which indicates that an increase in the riskiness of the financial
asset will reduce net saving; although the increase in risk diminishes
consumption, the decline in expected income from the decline in the
consumer's investment in the high-return asset is greater, and so
expected saving declines despite the decline in consumption.  It is
worth emphasizing again, though, that this decline in saving is {\it
  not} properly called a precautionary saving effect; it is a
consequence of the portfolio rebalancing which actually reduces the
size of the precautionary effect.  It would be confusing to refer to
this simply as a precautionary effect, when in fact the ultimate
outcome (increased saving) depends also on portfolio rebalancing,
income, and substitution effects.  This can be seen most clearly in
the case of logarithmic utility, where the precautionary effect is
precisely zero (see the discussion in \handoutM{CRRA-Portfolio}), yet
the portfolio and other effects exist and still generate the
conclusion that the increase in risk increases saving.

This analysis has been entirely in partial equilibrium.  General
equilibrium considerations (which arise, for example, in attempting to
use models of this kind to understand global imbalances) complicate
the picture further.  In particular, consumers' efforts to flee the
risky asset result in a lower equilibrium price for that asset (and a
correspondingly higher price for the riskless asset, and therefore a lower 
riskfree rate).  This induces
yet a {\it third} round of responses to the increase in risk.  

A final important observation is that the assumption that the
consumer has no labor income means that perhaps the largest
`classical' effect of interest rates on consumption and saving is
entirely omitted from the analysis here: The human wealth effect.
(See~\cite{summersCapTax} for a statement of the argument that the
human wealth effect of interest rates is likely in practice to be much
larger than the income and substitution effects.)  The general
equilibrium decline in riskfree rates attendant upon an increase in
riskiness of the risky asset should boost human wealth, increase
consumption, and reduce saving.  For careful and insightful treatments
of the general equilibrium problem, see \cite{paEntrep} and
\cite{corneliImbalances}.

\begin{figure}
\caption{$\MPC$ as a Function of $\sigma$ Calculated Using (Numerically) Exact and Approximate Methods} \label{fig:MPCvsSigma}
\includegraphics[width=6in]{../Figures/MPCApproxVsExactVsSigma}
\end{figure}

} % end ifExam

\end{verbatimwrite}
\input ./body.tex
\input handoutBibMake
\end{document}

